% ============================================================================
% Theorem 9': Sharp k-th Order LUT Interpolation Bounds
% (upgrades the DA optimality claim from cubic-only to all orders)
% ============================================================================

\subsection{Sharp k-th Order LUT Error Bounds}
\label{sec:lut-sharp}

The de~Boor bound $\varepsilon = M_2 h^2/8$ (Eq.~\ref{eq:de_boor_general})
governs \emph{linear} interpolation of $C^2$ activations. We now generalize
to LUTs of arbitrary interpolation order $k$ and establish the \emph{sharp
constants}---the best possible coefficients for every order---together with
an explicit extremal family attaining them. This upgrades the DA
optimality claim (Theorem~\ref{thm:da-optimality}) from an empirical
statement restricted to cubic polynomials to an exact, all-orders minimax
optimality theorem.

Define the \emph{window product constant}
\begin{equation}
Q_k \;=\; \max_{u \in [0, k-1]} \prod_{i=0}^{k-1} |u - i|,
\qquad k \geq 2.
\label{eq:window_product}
\end{equation}
For $k=2$ (linear interpolation), $Q_2 = 1/4$; for $k=3$ (quadratic
interpolation), $Q_3 = 1/(3\sqrt{3})$; for $k=4$, $Q_4 = 1$; for $k=5$,
$Q_5 \approx 3.631$. The \emph{sharp constant} is
$c_k = Q_k / k!$:
\begin{equation}
c_2 = \tfrac{1}{8}, \qquad
c_3 = \tfrac{1}{9\sqrt{3}} \approx 0.0642, \qquad
c_4 = \tfrac{1}{24} \approx 0.0417, \qquad
c_5 \approx 0.0303.
\label{eq:sharp_constants}
\end{equation}

\setcounter{theorem}{8}
\begin{theorem}[Sharp k-th Order LUT Error Bounds; Minimax Optimality of the DA Bound]
\label{thm:sharp-lut}
Let $\sigma \in C^k([a,b])$, $k \geq 2$, with $\|\sigma^{(k)}\|_\infty \leq M_k$.
Let $\mathrm{LUT}_k(\sigma)$ be the order-$k$ LUT on an equispaced grid of
cell width $h = (b-a)/N$---each window of $k$ consecutive nodes interpolated
by the degree-$(k-1)$ polynomial through them. Then:

\begin{enumerate}
    \item \textbf{Soundness (all orders, certified bound).}
    \begin{equation}
    \|\sigma - \mathrm{LUT}_k(\sigma)\|_\infty
    \;\leq\; \frac{M_k}{k!}\, Q_k\, h^k.
    \label{eq:sharp_soundness}
    \end{equation}

    \item \textbf{Exactness (sharpness).} The constant in
    Eq.~\eqref{eq:sharp_soundness} is \emph{attained exactly}: there exists
    $\sigma^{*} \in C^\infty$ with $\|\sigma^{*(k)}\|_\infty = M_k$ such that
    \begin{equation}
    \|\sigma^{*} - \mathrm{LUT}_k(\sigma^{*})\|_\infty
    \;=\; \frac{M_k}{k!}\, Q_k\, h^k,
    \label{eq:sharp_exactness}
    \end{equation}
    namely the extremal family
    $\sigma^{*}(x) = (M_k/k!) \prod_{i=0}^{k-1} (x - x_i)$ on each window
    (its $k$-th derivative is $M_k$ identically, and its LUT interpolant is
    the zero polynomial). Hence $\sup$ over the class equals the bound
    exactly---no closure or asymptotic argument is needed.

    \item \textbf{Minimax optimality of the abstract domain.} Among all
    sound affine (first-order) design-time schemes---any map from the
    $N{+}1$ LUT values to a piecewise-affine certificate---the worst-case
    certified bound over the class is at least
    $(M_k/k!)\,Q_k\,h^k$, and the LUT achieves it. For $k=2$ this
    recovers the classical optimal-recovery result: no affine method beats
    the de~Boor constant $M_2 h^2/8$.
\end{enumerate}
\end{theorem}

\vspace{2pt}
\noindent\textit{Proof sketch.}
\textbf{(1)} The pointwise interpolation remainder theorem gives, for
$\sigma \in C^k$ interpolated at nodes $x_0, \dots, x_{k-1}$,
\begin{equation}
|\sigma(x) - p(x)| = \frac{|\sigma^{(k)}(\xi_x)|}{k!}
\prod_{i=0}^{k-1} |x - x_i|
\;\leq\; \frac{M_k}{k!} \prod_{i=0}^{k-1} |x - x_i|.
\label{eq:interp_remainder}
\end{equation}
Writing $x = x_0 + u h$ with $u \in [0, k-1]$ gives
$\prod_i |x - x_i| = h^k \prod_i |u - i| \leq Q_k h^k$, whence
Eq.~\eqref{eq:sharp_soundness}. \textbf{(2)} For
$\sigma^{*}(x) = (M_k/k!)\prod_{i=0}^{k-1}(x - x_i)$, the leading
coefficient is $M_k/k!$, so $\sigma^{*(k)} \equiv M_k$; and
$\sigma^{*}$ vanishes at every node, so its LUT interpolant is the zero
polynomial. The error at $x$ is exactly
$(M_k/k!) \prod_i |x - x_i|$, whose window supremum is
$(M_k/k!) Q_k h^k$ by definition of $Q_k$
(Eq.~\eqref{eq:window_product}). \textbf{(3)} Any sound affine scheme
must be sound at $\sigma^{*}$ (a member of the class); since the true
supremum of $|\sigma^{*} - \mathrm{LUT}_k(\sigma^{*})|$ is the constant
above, any certificate bound below it is violated. The LUT attains it.
For $k = 2$ this is the classical minimax optimality of piecewise-linear
interpolation over the $C^2$ class. $\square$

\vspace{2pt}
\noindent\textbf{Remarks.}
\textit{(i) The folklore formula $M_k h^k / 2^k$ is false for $k \geq 3$.}
The sharp constants are $c_3 = 1/(9\sqrt{3}) \approx 0.0642 \neq 1/8$ and
$c_4 = 1/24 \neq 1/16$; the coincidental agreement holds only for $k \leq 2$.
\textit{(ii) The extremal family is a degree-$k$ polynomial with a specific
sign profile}---random cubic sampling (as in prior empirical claims) cannot
approach the constant, which is exactly why sharpness requires the explicit
construction above rather than Monte Carlo.
\textit{(iii) Compositional soundness.} For an $L$-layer SVNN with
activation layer $\ell$ having an order-$k_\ell$ LUT of width $h_\ell$,
\begin{equation}
B_{\mathrm{total}} \;=\;
\sum_{\ell=1}^{L} \Bigl( \prod_{j>\ell} \|W_j\|_\infty \Bigr)
\cdot \frac{M_k^{(\ell)}}{k_\ell!}\, Q_{k_\ell}\, h_\ell^{k_\ell}
\label{eq:sharp_compositional}
\end{equation}
is sound by induction (the pointwise error $e_\ell$ at layer $\ell$
satisfies $|e_\ell| \leq c_k M_k h^k$, and propagation through subsequent
linear maps amplifies by $\|W_j\|_\infty$). Each term is individually
attained by the sign-perturbed extremal
$\sigma_j = \pm \sigma^{*}$ (the involution $\sigma \mapsto -\sigma$
preserves $\|\sigma^{(k)}\|_\infty$). Exact \emph{simultaneous} attainment
holds for $L=1$; for $L \geq 2$ the bound is sharp up to the composition
(per-neuron sign freedom aligns intermediate pre-activations onto extremal
cells only under injectivity of the intermediate maps).

\vspace{2pt}
\noindent\textbf{Connection to DA (Theorem~\ref{thm:da-optimality}).}
Setting $k = 2$ in Theorem~\ref{thm:sharp-lut} yields the de~Boor bound
with sharp constant $1/8$ and the extremal quadratic
$\sigma^{*}(x) = (M_2/2)(x - x_0)(x - x_1)$, whose error attains
$M_2 h^2/8$ at the window midpoint---the semantic tightness of the DA
bound established in Theorem~\ref{thm:tightness} (King~C). The upgrade
here is threefold: (i)~the optimality claim extends from quadratic
functions (degree-$\leq 2$) to \emph{all} interpolation orders;
(ii)~the extremal family is explicit and verified to machine precision
($10{,}000$ random $C^k$ functions: soundness ratio $\leq 1$; extremal
ratio $= 1.000000$ exactly, E63); (iii)~the minimax optimality statement
is made formal over the class of affine design-time schemes.
